% macros.tex — symbol macros + theorem-environment setup.
% Generated for the four Lean-verified results in `.lean-proof/Proof/Proofs/DecodeIff.lean`
% (`Decode.decode_iff_gradient_separation`, `..._descent_separation_of_neg`,
% `..._fails_iff_some_competitor`, `..._softmax_residual`). Symbol macros are transcribed
% verbatim from `.lean-proof/notation-map.md`; the theorem-environment block follows the
% cleveref-safe \newaliascnt template in references/conventions.md (lint R0c).

% ---------------------------------------------------------------------------
% Symbol macros (one per row of the notation map)
% ---------------------------------------------------------------------------
\newcommand{\Vocab}{\mathcal{V}}                                  % finite vocabulary
\newcommand{\inner}[2]{\left\langle #1, #2 \right\rangle}         % real inner product
\newcommand{\norm}[1]{\left\| #1 \right\|}                         % norm
\newcommand{\Sphere}{\mathcal{S}_r}                                % LayerNorm sphere of radius r
\newcommand{\loss}{L}                                              % implicit loss
\newcommand{\grad}{\nabla}                                         % gradient operator
\newcommand{\xstar}{x^{\star}}                                     % stationary point
\newcommand{\astar}{a^{\star}}                                     % designated (answer) token
\newcommand{\Generated}{\mathsf{Gen}}                              % decoder predicate
\newcommand{\indicator}[1]{\mathbf{1}\!\left[#1\right]}            % one-hot target indicator 1[c = a*]

% --- 2nd-layer (bridge + dynamics, sections/04) symbols ---
\newcommand{\Zpart}{Z}                                            % softmax partition Z = sum_c exp<W_c,x>
\newcommand{\logit}[1]{\ell_{#1}}                                  % logit ell_c = <W_c, x>
\newcommand{\optgap}{\Delta}                                      % suboptimality margin Delta = log 2 - L*
\newcommand{\lossopt}{\loss^{\star}}                              % optimal loss value L*
\newcommand{\tstar}{t^{\star}}                                    % decode-time threshold t*
\newcommand{\passat}[1]{\mathrm{pass}@#1}                          % pass@T success rate over a problem distribution
\newcommand{\smooth}{\beta}                                       % smoothness constant beta
\newcommand{\stepsize}{\eta}                                      % GD step size eta

% ---------------------------------------------------------------------------
% Theorem-environment setup (aliascnt → cleveref-safe shared numbering)
% ---------------------------------------------------------------------------
\usepackage{amsthm}
\usepackage{aliascnt}
\usepackage[capitalize]{cleveref}

% Master counter.
\theoremstyle{plain}
\newtheorem{theorem}{Theorem}[section]

% Aliased plain-style environments.
\newaliascnt{lemma}{theorem}
\newtheorem{lemma}[lemma]{Lemma}
\aliascntresetthe{lemma}

\newaliascnt{corollary}{theorem}
\newtheorem{corollary}[corollary]{Corollary}
\aliascntresetthe{corollary}

\newaliascnt{proposition}{theorem}
\newtheorem{proposition}[proposition]{Proposition}
\aliascntresetthe{proposition}

% Aliased definition-style environments.
\theoremstyle{definition}
\newaliascnt{definition}{theorem}
\newtheorem{definition}[definition]{Definition}
\aliascntresetthe{definition}

% Aliased remark-style environment.
\theoremstyle{remark}
\newaliascnt{remark}{theorem}
\newtheorem{remark}[remark]{Remark}
\aliascntresetthe{remark}

% Cleveref names — keyed on the aliascnt counters, so they fire correctly.
\crefname{theorem}{Theorem}{Theorems}
\crefname{lemma}{Lemma}{Lemmas}
\crefname{corollary}{Corollary}{Corollaries}
\crefname{proposition}{Proposition}{Propositions}
\crefname{definition}{Definition}{Definitions}
\crefname{remark}{Remark}{Remarks}
\crefname{equation}{Eq.}{Eqs.}
